\documentclass[aps,prx,reprint,superscriptaddress,nofootinbib,floatfix]{revtex4-2}

\usepackage{graphicx}
\usepackage{amsmath,amssymb}
\usepackage{bm}
\usepackage{xcolor}
\usepackage[colorlinks=true,allcolors=blue]{hyperref}
\usepackage{float}

\graphicspath{{figures/}{../figures/}}

\renewcommand{\thesection}{S\arabic{section}}
\renewcommand{\theequation}{S\arabic{equation}}
\renewcommand{\thetable}{S\arabic{table}}
\renewcommand{\thefigure}{S\arabic{figure}}

\begin{document}

\title{Supplemental Material:\\
What a radical-pair quantum reservoir would require}

\author{Hikaru Wakaura}
\affiliation{QIRI (Quantum Integrated Research Institute Inc.), Tokyo, Japan}
\author{Taiki Tanimae}
\affiliation{QIRI (Quantum Integrated Research Institute Inc.), Tokyo, Japan}

\date{\today}

\maketitle

Every number below is computed under the single protocol used in the main text ---
correlated-pair birth, converged time grid, out-of-sample capacity estimation ---
so that no quantity is carried over from a superseded calculation. The sections cover
(S1) validation of the Liouville-space solver, (S2) the capacity estimator and its
controls, (S3) the injection map and why a single-electron reset fails, (S4) the
chemically accessible readout routes and their grid convergence, (S5) the coherence
fraction at matched hopping rate, (S6) classical baselines, (S7) the turnover clock
and the register-reuse criterion, and (S8) the spin parameters. All are reproducible
from the open \texttt{qbscreen} package.

\section{Liouville-space solver and the coherent limit}
\label{sec:solver}
The propagator is $e^{\mathcal{L}\tau}$ of the Haberkorn master equation assembled in
Liouville space with the column-stacking convention
$\mathrm{vec}(A\rho B)=(B^{\!\top}\!\otimes A)\,\mathrm{vec}(\rho)$; the
time-integrated singlet yield follows from a single linear solve,
\begin{equation}
\Phi_S = k_S\,\mathrm{vec}(P_S)^{\dagger}\,(-\mathcal{L})^{-1}\,\mathrm{vec}(\rho_0).
\label{eq:yield}
\end{equation}
Adding the decoherence term must not corrupt the coherent dynamics. For
$1/T_2^e\to0$ with $k_S=k_T=k$, Eq.~\eqref{eq:yield} must reproduce the closed-form
eigenbasis result
$\Phi_S^{\mathrm{coh}}=k\big[\sum_m P_{mm}\rho_{mm}/k+\mathrm{Re}\sum_{m\neq n}
P_{mn}\rho_{nm}/(k+i\omega_{mn})\big]$. For a two-electron-plus-nucleus model at
$B=0.9$~mT, $A=200$~MHz and $k=1~\mu\mathrm{s}^{-1}$ the two agree to better than
$10^{-6}$; this is asserted as a unit test.

\section{Capacity estimator and its controls}
\label{sec:capacity}
For a target $y_k$ of the input history the capacity is the out-of-sample squared
multiple correlation: a ridge readout ($\lambda=10^{-6}$) is fitted on the first half
of the run and scored on the held-out second half. The totals are
$\mathrm{MC}=\sum_{d\ge0}C[s_{k-d}]$ and
$\mathrm{IPC}=\mathrm{MC}+\sum_{\rm nonlin}C[y]$ over the degree-2 Legendre family
(diagonal targets to $d_{\max}=14$, cross targets for $d_1<d_2\le6$).

\paragraph*{Null floor.} Scoring the same feature matrices against targets built from
an \emph{independent} input sequence, for which the true capacity is zero, returns
$+0.010$ for the estimator used here. The same sum without the $\max(0,\cdot)$ clip
returns $-1.49$, i.e.\ it is far more biased, in the opposite direction. We therefore
retain the clip and quote the floor rather than removing it.

\paragraph*{Convergence.} The capacity is converged in the length $L$ of the driving
sequence (six seeds, engineered point): $\mathrm{IPC}=5.28\pm0.26$, $5.60\pm0.20$,
$5.87\pm0.19$ and $5.84\pm0.10$ for $L=600$, $1200$, $2400$ and $4800$, i.e.\ flat
beyond $L\approx2400$. Under the superseded single-electron-reset protocol the same
estimator did \emph{not} converge, rising from $5.9$ to $8.1$ over the same range;
convergence is a property of the corrected drive, not of the estimator.

\paragraph*{The Dambre bound is not an independence test.} $\mathrm{IPC}\le
n_{\rm obs}$ holds asymptotically for an orthonormal target family but is violated in
finite samples: for $X=s$ ($n_{\rm obs}=1$, $T=1200$) the estimator returns $1.011$
on average and exceeds the bound in $30$ of $30$ seeds. Satisfying the inequality
therefore does not demonstrate that the readout observables are independent, and we
make no such argument anywhere in this work.

\section{The injection map}
\label{sec:injection}
A protocol that resets only the input electron to a product state while retaining the
reduced state of the rest is not a chemical process. For a separated pair ($J=0$) it
destroys the electron--electron correlation on every cycle, and since
$\langle P_S\rangle=\tfrac14-\langle \bm S_1\!\cdot\!\bm S_2\rangle$ with
$\langle \bm S_2\rangle=0$, the singlet probability is pinned at exactly $1/4$
independently of the input. We replace it by correlated-pair birth: the nuclei are
retained and the electron pair is recreated as
$\rho_e(s)=s\,|S\rangle\langle S|+(1-s)P_T/3$, or as a coherent $S$--$T_0$
superposition. Table~\ref{tab:inject} gives the consequence at the cryptochrome point
($B=50~\mu$T, $J=0$, $\tau=1~\mu$s, six seeds, kinetic-only readout).

\begin{table}[H]
\caption{\label{tab:inject}The injection map and the capacity it permits.}
\begin{ruledtabular}
\begin{tabular}{lcc}
injection & max.\ flux spread & IPC \\
\hline
single-electron reset (superseded) & $2.7\times10^{-14}$ & $0.000\pm0.000$ \\
correlated-pair birth, $S/T$ & $4.9\times10^{-1}$ & $2.941\pm0.108$ \\
correlated-pair birth, $S$--$T_0$ & $4.3\times10^{-1}$ & $3.116\pm0.209$ \\
\end{tabular}
\end{ruledtabular}
\end{table}

Under the superseded map the spread of the observable is at the level of machine
precision: the readout is a constant, and the vanishing capacity is an identity of
the map rather than a statement about spin chemistry.

\section{Chemically accessible readout routes}
\label{sec:routes}
Recombination is switched on ($k_S=1$, $k_T=0.2~\mu\mathrm{s}^{-1}$) so that the
observable is the yield actually formed,
$Y_X=k_X\!\int_0^\tau\!\mathrm{Tr}[P_X\rho(t)]\,\mathrm{d}t$, sampled at five
laboratory times. CIDNP is the product-weighted nuclear polarisation,
$\langle I_z\rangle_{\rm prod}=\int k_S\mathrm{Tr}[I_zP_S\rho]\,\mathrm{d}t\,/
\int k_S\mathrm{Tr}[P_S\rho]\,\mathrm{d}t$: the polarisation carried away by the
recombined molecule, not that of the surviving pair.

\begin{table}[H]
\caption{\label{tab:routes}Capacity recovered by chemically accessible observables
(six seeds).}
\begin{ruledtabular}
\begin{tabular}{lcc}
readout route & channels & IPC \\
\hline
end-point singlet yield & 1 & $1.017$ \\
both exit channels, end point & 2 & $1.374$ \\
three micro-environments (different $\tau$) & 3 & $2.557$ \\
time-resolved kinetics, one product pool & 5 & $2.008$ \\
\quad $+$ accumulated downstream pool & 10 & $2.706$ \\
\quad $+$ CIDNP on the product & 8 & $4.650$ \\
\end{tabular}
\end{ruledtabular}
\end{table}

\paragraph*{Time-grid convergence.} A $40$~MHz hyperfine coupling has a $25$~ns
period, so a coarse grid under-samples the dynamics and inflates the time-resolved
capacities. Table~\ref{tab:grid} shows convergence from $n_t=48$ onward; the main
text uses $n_t=96$ ($\mathrm{d}t\approx10$~ns).

\begin{table}[H]
\caption{\label{tab:grid}Convergence in the number of time steps per cycle (six seeds, $L=700$, the same protocol as Table~\ref{tab:routes}).}
\begin{ruledtabular}
\begin{tabular}{lccc}
$n_t$ & $\mathrm{d}t$ (ns) & kinetics IPC & $+$CIDNP IPC \\
\hline
24 & 41.7 & 2.756 & 4.879 \\
48 & 20.8 & 2.050 & 4.614 \\
96 & 10.4 & 2.008 & 4.650 \\
192 & 5.2 & 1.988 & 4.656 \\
\end{tabular}
\end{ruledtabular}
\end{table}

\paragraph*{Why CIDNP is not an optional extra channel.} The reduced state of either
electron of a newly born pair is maximally mixed for every input, because both
$|S\rangle\langle S|$ and $P_T/3$ carry zero one-electron polarisation; the input
resides entirely in the two-electron correlation. With $J=0$ the radicals do not
interact, so nothing transfers that correlation to a nucleus. Consistently, with
recombination switched off the cryptochrome-point capacity is $1.01$, the memoryless
floor. Spin-selective recombination is what converts the singlet--triplet correlation
into a population difference and so writes the input into the nuclear register: in a
separated pair it is the write mechanism, not merely the read mechanism.

\section{Coherence fraction at matched hopping rate}
\label{sec:coh}
An all-spin pure-dephasing rate $\gamma$ is added in the $S_z$ product basis. Taking
$\gamma\to\infty$ at fixed $\tau$ measures Zeno freezing rather than a classical
limit, because the golden-rule transfer rate scales as $|H_{ij}|^2/\gamma$; we
therefore hold rate\,$\times\,\tau$ fixed by scaling $\tau\propto\gamma$, whereupon
the capacity saturates.

\begin{table}[H]
\caption{\label{tab:coh}Incoherent limit at matched hopping rate (six seeds).}
\begin{ruledtabular}
\begin{tabular}{lcc}
$\gamma$ (MHz) & $\tau$ (ns) & IPC \\
\hline
0 (coherent) & 50 & $5.421$ \\
$10^{3}$ & $10^{3}$ & $0.999$ \\
$5\times10^{3}$ & $5\times10^{3}$ & $0.998$ \\
$2\times10^{4}$ & $2\times10^{4}$ & $0.998$ \\
\end{tabular}
\end{ruledtabular}
\end{table}

The saturation value $1.00$ is the memoryless floor, so the coherence fraction is
$(5.421-0.999)/5.421=0.82$.

\section{Classical baselines}
\label{sec:esn}
Tanh echo-state networks (spectral radius $0.9$, input scale $0.6$) are driven by the
same input and scored with the same estimator (eight seeds, $L=1200$).

\begin{table}[H]
\caption{\label{tab:esn}The quantum reservoir against matched classical reservoirs.}
\begin{ruledtabular}
\begin{tabular}{lc}
reservoir (readout features) & IPC \\
\hline
quantum, 5 spins (12 observables) & $5.593\pm0.191$ \\
ESN, 5 nodes (5 linear) & $4.552\pm0.325$ \\
ESN, 5 nodes (12: 5 linear $+$ 7 quadratic) & $9.414\pm0.860$ \\
ESN, 12 nodes (12 linear) & $10.366\pm0.590$ \\
\end{tabular}
\end{ruledtabular}
\end{table}

Matched by physical unit the quantum system is ahead; matched by readout dimension it
is behind. We report both and claim no advantage. Which accounting is fair depends on
a cost model for non-commuting measurements against multiplications that we do not
have.

\section{Turnover clock and the register-reuse criterion}
\label{sec:clock}
The memory horizon is $\mathrm{MC}$ times the cycle time. Because the nuclei leave in
the diamagnetic product, where nuclear relaxation is of order $0.1$--$1$~s, the cycle
time is the turnover interval $T_{\rm d}$ and not the radical-pair lifetime. We
insert a pause between turnovers and relax the register with probability
$q=1-\exp(-T_{\rm d}/T_{1n})$, taking $T_{1n}=1$~s.

\begin{table}[H]
\caption{\label{tab:clock}The horizon follows the turnover interval.}
\begin{ruledtabular}
\begin{tabular}{lcccc}
$T_{\rm d}$ & $q$ & MC (5 ch) & MC (8 ch) & horizon (8 ch) \\
\hline
$1~\mu$s & $10^{-6}$ & $1.948$ & $2.818$ & $5.6~\mu$s \\
$1$~ms & $10^{-3}$ & $1.948$ & $2.818$ & $2.8$~ms \\
$10$~ms & $10^{-2}$ & $1.947$ & $2.817$ & $28$~ms \\
$100$~ms & $9.5\times10^{-2}$ & $1.944$ & $2.800$ & $280$~ms \\
\end{tabular}
\end{ruledtabular}
\end{table}

The capacity changes by $0.2\%$ across five decades of turnover interval, so the
horizon simply tracks it. Table~\ref{tab:reuse} isolates the requirement that
actually binds: how much depolarisation of the register between turnovers the memory
tolerates.

\begin{table}[H]
\caption{\label{tab:reuse}The criterion: memory against register depolarisation $q$.}
\begin{ruledtabular}
\begin{tabular}{lccc}
$q$ & $T_{\rm d}/T_{1n}$ & MC (5 ch) & MC (8 ch) \\
\hline
0.00 & 0 & $1.984$ & $2.890$ \\
0.30 & 0.36 & $1.958$ & $2.667$ \\
0.50 & 0.69 & $1.927$ & $2.288$ \\
0.70 & 1.20 & $1.788$ & $2.030$ \\
0.90 & 2.30 & $1.218$ & $1.977$ \\
1.00 & $\infty$ & $1.002$ & $1.002$ \\
\end{tabular}
\end{ruledtabular}
\end{table}

Memory is essentially intact while the turnover interval is shorter than the nuclear
relaxation time ($q\lesssim0.5$) and collapses to the memoryless floor
$\mathrm{MC}=1.00$ when every turnover starts from a fresh molecule.

\section{Spin parameters}
\label{sec:params}
Each radical carries its own nuclei; they are not shared, which is what a spatially
separated pair means. The Hamiltonian therefore contains three hyperfine terms, two
on the first radical and one on the second. An earlier version of the code accepted a
fourth hyperfine argument that was never added to the Hamiltonian and had no effect on
any result; it has been removed rather than left in place.

\begin{table}[H]
\caption{\label{tab:params}Hyperfine assignment at the cryptochrome point. Values are
representative isotropic magnitudes of the dominant FAD$^{\bullet-}$ and
TrpH$^{\bullet+}$ couplings; the full anisotropic tensors are given by Lee
\textit{et al.}, J.~R.~Soc.~Interface \textbf{11}, 20131063 (2014), and the
flavin-radical ENDOR couplings by Schleicher \textit{et al.}, Sci.~Rep.\
\textbf{11}, 18395 (2021).}
\begin{ruledtabular}
\begin{tabular}{lccc}
coupling & radical / nucleus & (mT) & (MHz) \\
\hline
$A_{e_1a}$ & FAD$^{\bullet-}$, $^{14}$N5 & $\sim0.50$ & $14$ \\
$A_{e_1b}$ & FAD$^{\bullet-}$, $^{14}$N10 & $\sim0.18$ & $5$ \\
$A_{e_2a}$ & TrpH$^{\bullet+}$, H$\beta$1 & $\sim1.43$ & $40$ \\
\end{tabular}
\end{ruledtabular}
\end{table}

The engineered operating point uses $A_{e_1a}=80$, $A_{e_1b}=40$, $A_{e_2a}=15$~MHz,
$J=2$~MHz, $B=1$~mT and $\tau=T_2^e=50$~ns. The conclusions depend on these
magnitudes lying in the few--$40$~MHz range typical of flavin and tryptophan
radicals, not on the precise values; the register criterion of Sec.~\ref{sec:clock}
is independent of the hyperfine assignment altogether.

\end{document}
